\section{Связность графов. Алгоритмы обхода графов}
\label{sec:graph-connectivity}

\subsection{Компоненты связности}

Две вершины неориентированного графа называются связными, если существует соединяющая их цепь. Отношение связности является отношением эквивалентности, поэтому множество вершин графа разбивается на \textbf{компоненты связности}.

Граф называется \textbf{связным}, если он имеет ровно одну компоненту связности.

Ребро называется \textbf{мостом}, если его удаление увеличивает число компонент связности. Вершина называется \textbf{точкой сочленения}, если её удаление вместе с инцидентными ей рёбрами увеличивает число компонент связности.

Для связного графа ребро является мостом тогда и только тогда, когда оно не принадлежит ни одному циклу.

Для количественной характеристики связности вводятся следующие величины:
\begin{itemize}
	\item \textbf{вершинная связность} $\kappa(G)$ --- минимальное число вершин, удаление которых делает граф несвязным или тривиальным;
	\item \textbf{рёберная связность} $\lambda(G)$ --- минимальное число рёбер, удаление которых нарушает связность графа.
\end{itemize}

Для нетривиального графа выполняется неравенство
\[
\kappa(G)\leq\lambda(G)\leq\delta(G),
\]
где $\delta(G)$ --- минимальная степень вершины.

Если связный граф имеет точку сочленения, то $\kappa(G)=1$. Если связный граф имеет мост, то $\lambda(G)=1$.

\subsection{Теорема Менгера}

Два $u$--$v$-пути называются \textbf{внутренне вершинно непересекающимися}, если у них нет общих вершин, кроме $u$ и $v$. Пути называются \textbf{рёберно непересекающимися}, если они не имеют общих рёбер.

Внутренне вершинно непересекающиеся $u$--$v$-пути являются также рёберно непересекающимися.

\textbf{Теорема Менгера, вершинная форма.}
Пусть $u$ и $v$ --- две несмежные вершины графа. Тогда минимальное число вершин, удаление которых разделяет $u$ и $v$, равно максимальному числу попарно внутренне вершинно непересекающихся
$u$--$v$-путей:
\[
\max |P(u,v)|=\min |S(u,v)|.
\]

\textbf{Теорема Менгера, рёберная форма.}
Минимальное число рёбер, удаление которых разделяет вершины $u$ и $v$, равно максимальному числу попарно рёберно непересекающихся
$u$--$v$-путей.

В частности, граф является $k$-связным тогда и только тогда, когда любые две несмежные вершины соединены не менее чем $k$ попарно внутренне вершинно непересекающимися путями.

Теорема Менгера связывает структурную устойчивость графа с числом независимых маршрутов между его вершинами. Она также тесно связана с теоремой о максимальном потоке и минимальном разрезе.

\subsection{Поиск в ширину}

\textbf{Поиск в ширину} (\textit{Breadth-First Search}, BFS) обходит граф слоями по расстоянию от начальной вершины. Для организации обхода используется очередь.

Алгоритм:
\begin{enumerate}
	\item отметить стартовую вершину $s$, положить $dist[s]=0$ и поместить её в очередь;
	\item пока очередь не пуста, извлечь из неё вершину $v$;
	\item для каждого неотмеченного соседа $u$ вершины $v$:
	\begin{itemize}
		\item отметить $u$;
		\item положить $parent[u]=v$;
		\item положить $dist[u]=dist[v]+1$;
		\item добавить $u$ в очередь.
	\end{itemize}
\end{enumerate}

При представлении графа списками смежности каждая вершина посещается не более одного раза, а каждое ребро просматривается постоянное число раз, поэтому
\[
T(n,m)=O(n+m).
\]

Дополнительная память алгоритма составляет $O(n)$. Если учитывать также хранение графа в списках смежности, требуется $O(n+m)$ памяти.

BFS вычисляет расстояния от стартовой вершины до всех достижимых вершин невзвешенного графа, где длина пути измеряется числом рёбер. Массив $parent$ задаёт дерево кратчайших путей из стартовой вершины.

\subsection{Поиск в глубину}

\textbf{Поиск в глубину} (\textit{Depth-First Search}, DFS) проходит из текущей вершины в ещё не посещённую соседнюю вершину и продолжает движение максимально глубоко, после чего возвращается назад.

DFS может быть реализован рекурсивно или с помощью явного стека.

Базовая рекурсивная схема:
\begin{enumerate}
	\item отметить текущую вершину $v$ как посещённую;
	\item для каждого соседа $u$ вершины $v$, если $u$ ещё не посещена, вызвать DFS$(u)$.
\end{enumerate}

При использовании списков смежности время работы равно
\[
O(n+m).
\]

Дополнительная память составляет $O(n)$ и требуется для массива посещённости, служебных массивов и стека DFS.

Если граф несвязен, внешний цикл запускает DFS из каждой ещё не посещённой вершины. В результате получается \textbf{лес DFS}, содержащий по одному дереву на каждую компоненту связности.

DFS применяется для:
\begin{itemize}
	\item поиска компонент связности;
	\item обнаружения циклов;
	\item поиска мостов и точек сочленения;
	\item топологической сортировки;
	\item поиска компонент сильной связности.
\end{itemize}


\subsection{Связность ориентированных графов}

Орграф называется \textbf{слабо связным}, если после игнорирования направлений всех дуг получается связный неориентированный граф.

Орграф называется \textbf{сильно связным}, если для любых двух вершин $u,v$ существуют ориентированные пути
\[
u\leadsto v
\qquad\text{и}\qquad
v\leadsto u.
\]

\textbf{Компонента сильной связности} --- максимальное по включению множество вершин, попарно взаимно достижимых друг из друга.

После стягивания каждой компоненты сильной связности в одну вершину получается ориентированный ациклический граф, называемый \textbf{конденсацией} графа.

Для поиска компонент сильной связности можно использовать \textbf{алгоритм Косарайю}:
\begin{enumerate}
	\item выполнить DFS по исходному графу $G$ и запомнить вершины в порядке завершения их обработки;
	\item построить транспонированный граф $G^T$, в котором направления всех дуг обращены;
	\item запускать DFS по $G^T$ (в каждой непосещённой вершине) в порядке убывания времени выхода вершин. Множество вершин, найденное каждым запуском DFS, образует одну компоненту сильной связности.
\end{enumerate}

Каждый обход DFS и построение транспонированного графа выполняются за $O(n+m)$, поэтому общая сложность алгоритма равна 
\[
O(n+m).
\]

\subsection{Топологическая сортировка}

\textbf{Топологическим порядком} ориентированного графа называется такой линейный порядок его вершин, в котором для каждой дуги $(u,v)$ вершина $u$ расположена раньше вершины $v$.

Топологический порядок существует тогда и только тогда, когда ориентированный граф не содержит ориентированных циклов, то есть является DAG (\textit{Directed Acyclic Graph}).

Топологическую сортировку можно получить DFS. Вершина добавляется в список после завершения обработки всех исходящих из неё рёбер. Обратный порядок выхода из вершин является топологическим порядком, если граф ацикличен. При DFS наличие дуги в вершину, находящуюся в состоянии обработки, свидетельствует о наличии ориентированного цикла.

Алгоритм при использовании списков смежности работает за
\[
O(n+m).
\]

\subsection{Что важно сказать на экзамене}

BFS и DFS --- базовые алгоритмы обхода графов с линейной сложностью
\[
O(n+m).
\]

BFS обходит граф слоями и находит кратчайшие пути по числу рёбер в невзвешенном графе.

DFS строит дерево или лес поиска и позволяет исследовать структурные свойства графа: компоненты сильной связности и выполнять топологическую сортировку.

Теорема Менгера связывает связность графа с числом независимых путей:
минимальное число разделяющих вершин или рёбер совпадает с максимальным
числом соответствующих непересекающихся путей.

\newpage